\def\oddtext{People have \hbox{$\alpha\beta\delta\gamma\epsilon221$} commented that there are ridiculous conferences calculating expressions like \hbox{$1+ 2+ 6 \alpha\beta+2x ( y+z ) dx$}.}

This is a quick plain-TeX highlighting just a small number of errors and warnings.
The work I've done so far on the ``expressive'' diagnostics is just scratching the surface; there are a lot of errors that can crop up in TeX. 
Some can be helped with almost no changes to the underlying engine; some would require major changes (like trying to track the provenance of tokens).
Doing those changes in either a performance-neutral (or, where not possible, opt-in) way is on the roadmap.

Lot's of errors are typos:
$\vce{a} \cdoot \ahat{n} = 12$

Giving better errors for these can be done with the information the TeX82 engine already has.
What about times when we want to report the typeset lines rather than the source... after all, the lines might have come from anywhere

\hsize=2.33in %force our lines to be quite small!

Several ridiculous conferences routinely characterize frivolous phenomenology.

\hsize=2.5in

\oddtext\par

\hsize=6in

Another kind of error are those pesky ``oh, you forgot to close your math environment'' variety:

I have computed this expression: $10 = 1+2+3+4

We forgot to close the math environment there... whoops!

Having diagnostics pointing not only to the insertion point but also where the mode was opened requires a little bit of extra information to be stored (in the list\_state\_record, if you're interested), but not much.
Some examples:
$ 10 = 1 + 2
   +3+4

Or what if we start and end in different files?
$ 10 = 1 + \input error_math_stuff

\openmath{x^2 + y^2} = r^2 

This conveniently highlights the interaction between the semantic nest, the scanner, etc.
With more work we could decide to highlight the macro, or move where that down-pointing ``caret'' is located, etc.


\bye
